% !TEX root = ../notes.tex

\documentclass[../notes.tex]{subfiles}

\begin{document}

\section{March 11}
We start class by saying a little about the Langlands dual.

\subsection{The Langlands Dual}
Let's attempt to generalize our construction of $\widehat T$. Recall that split reductive groups are classified by their root data.
\begin{definition}[root datum]
	Fix a connected split reductive group $G$ over a field $F$, and choose a split maximal torus $T\subseteq G$. Then we define its \textit{root datum} as the quadruple $(\Lambda,R,\Lambda^\lor,R^\lor)$, where $\Lambda=X^*(T)$ is the weight lattice, $\Lambda^\lor=X_*(T)$ is the coweight lattice, $R$ consists of the roots, and $R^\lor$ consists of the coroots.
\end{definition}
\begin{remark}
	Semisimple groups are classified merely by the root system $(\Lambda,R)$. Remembering the coroots is required to access all split reductive groups.
\end{remark}
Here is the main theorem.
\begin{theorem} \label{thm:root-datum}
	Isomorphism classes of connected split reductive groups are equivalent to isomorphism classes of root data.
\end{theorem}
\begin{proof}
	This is similar to the argument for connected semisimple groups, but now we attempt to remember the central torus of the group. We remark that the existence of this group is difficult.
\end{proof}
We are thus allowed to make the following definition.
\begin{definition}[Langlands dual]
	Fix a connected split reductive group $G$ over a field $F$ with root datum $(\Lambda,R,\Lambda^\lor,R^\lor)$. Then the \textit{Langlands dual} $\widehat G$ is the complex connected split reductive group with root datum $(\Lambda^\lor,R^\lor,\Lambda,R)$.
\end{definition}
\begin{remark}
	We can see that the type of the dual group is the one given by reversing the Dynkin diagram.
\end{remark}
\begin{example}
	Take $G\coloneqq\op{GL}_n$, for which we can take the diagonal torus. Then $X^*(T)$ and $X_*(T)$ both have a canonical basis $\{e_i\}$ (which is self-dual). The root system of $\op{GL}_n$ is $\{e_i-e_j\}$, and the coroots are given by $\{e_i^*-e_j^*\}$, so we conclude that the Langlands dual group is $\op{GL}_{n,\CC}$.
\end{example}
\begin{example}
	Take $G\coloneqq\op{SL}_n$, for which we can again take the diagonal torus. Then $X^*(T)=\ZZ^n/\Delta\ZZ$ (meaning that we are taking a quotient by the ``diagonal'' copy of $\ZZ$, arising because determinants must vanish). Then the root system is still $\{e_i-e_j\}$ for the standard basis $\{e_i\}$ of $X^*(T)$. On the other hand, $X_*(T)$ is the trace-zero hyperplane $\ZZ^n_0$ of $\ZZ^n$, and the coroot system is again just given by $\{e_i^*-e_j^*\}$. Redoing the argument shows that the dual root data is given by $\op{PGL}_n$, so $\widehat G=\op{PGL}_{n,\CC}$.
\end{example}
\begin{example}
	Take $G=\op{Sp}_{2n}$. Using the standard symplectic form, one finds that $X^*(T)=\ZZ^n$ with dual $X_*(T)=\ZZ^n$. The root system of $\op{Sp}_{2n}$ is given by $\{e_i-e_j\}_{i\ne j}\sqcup\{2e_i\}$, which only has the long root $2e_i$. On the other hand, the coroot system is given by $\{e_i^*-e_j^*\}_{i\ne j}\sqcup\{e_i^*\}$, which has only the short root $e_i^*$. Redoing the argument shows that the dual root data is given by $\op{SO}_{2n+1}$, so $\widehat G=\op{SO}_{2n+1,\CC}$.
\end{example}
The moral is that we are able to restate our classification of spherical representations as follows.
\begin{theorem}
	Fix a split connected reductive group $G$ over a nonarchimedean local field $F$. Then the spherical representations of $G(F)$ are in bijection with the semisimple conjugacy classes in $\widehat G(\CC)$.
\end{theorem}
\begin{proof}
	In light of our prior discussion, it only remains to show that the semisimple conjugacy classes in $\widehat G(\CC)$ are in bijection with $\widehat T/W$. Well, any semisimple element has a conjugate in $\widehat T$, which is unique up action by the Weyl group $W=\op N_G(T)/T$.
\end{proof}
\begin{remark}
	To explain how this generalizes, note that we can think of semisimple conjugacy classes as homomorphisms $W_F\to\widehat G(\CC)$ (up to conjugacy) which factor through $\ZZ$ and have semisimple image. It is a major area of research to figure out how to parameterize arbitrary representations of $G(F)$ by arbitrary homomorphisms $W_F\to\widehat G(\CC)$.
\end{remark}
\begin{remark}
	There is an analogous statement for real groups. The correct definition of Weil group $W_\RR$ is fitting into a non-split short exact sequence
	\[1\to\CC^\times\to W_\RR\to\op{Gal}(\CC/\RR)\to1.\]
\end{remark}

\end{document}